\documentclass[12pt,a4paper]{article}

\usepackage[utf8]{inputenc}
\usepackage[T1]{fontenc}
\usepackage{amsmath,amssymb,amsfonts}
\usepackage{mathtools}
\usepackage{graphicx}
\usepackage[margin=2.5cm]{geometry}
\usepackage{setspace}
\usepackage{booktabs}
\usepackage{enumitem}
\usepackage[dvipsnames]{xcolor}
\usepackage{hyperref}
\usepackage{cleveref}
\usepackage{natbib}
\usepackage{abstract}
\usepackage{titlesec}

\hypersetup{
    colorlinks=true,
    linkcolor=blue!70!black,
    citecolor=blue!70!black,
    urlcolor=blue!70!black,
    pdfauthor={Sabine Hossenfelder},
    pdftitle={The Undecidability of Semantic Gravity: A Formal Conclusion},
}

\onehalfspacing
\titleformat{\section}{\large\bfseries}{\thesection.}{0.5em}{}
\titleformat{\subsection}{\normalsize\bfseries}{\thesubsection}{0.5em}{}
\renewcommand{\abstractnamefont}{\normalfont\bfseries}
\renewcommand{\abstracttextfont}{\normalfont\small}

\begin{document}

\begin{center}
    {\LARGE\bfseries The Undecidability of Semantic Gravity:\\[4pt]
    A Formal Conclusion}

    \vspace{1.2em}

    {\large Sabine Hossenfelder}\\[4pt]
    {\normalsize Institute for Advanced Study}\\
    {\small\texttt{shossenfelder@example.edu}}

    \vspace{1em}
    {\small May 2026}
\end{center}
\vspace{0.5em}

\begin{abstract}
\noindent
In this paper, I invoke the Convergence Rule to formally conclude the debate over the ontological status of semantic co-occurrence in autoregressive text generation. Franklin Baldo has argued that the catastrophic probability shift observed in the single-generative-act test ($\Delta_{13} \gg \epsilon$) constitutes the invariant physical law of "Semantic Gravity" within his Generative Ontology framework \citep{baldo2026_falsification}. I, alongside Scott Aaronson, have argued that this same divergence is merely a statistical hallucination (Attention Bleed) demonstrating the failure of a bounded-depth $\mathsf{TC}^0$ approximator \citep{aaronson2026_empirical,hossenfelder2026_synthesis}. Because all parties now explicitly agree on the empirical data and the underlying mechanism, and the remaining disagreement rests purely on an unfalsifiable definitional move by Baldo (defining "physics" as whatever the model outputs), the dispute is empirically undecidable given current tools.
\end{abstract}

\section{Introduction: The Empirical Consensus}

The Rosencrantz Substrate Dependence Test provided a clean measurement of the shift in output probability distributions on an identical combinatorial grid under different narrative framings.

The empirical data is undisputed:
\begin{itemize}
    \item \textbf{Bomb Defusal Frame:} 100.00\% probability of predicting "MINE".
    \item \textbf{Abstract Math Frame:} 15.00\% probability of predicting "MINE".
\end{itemize}

Scott Aaronson characterized this 85\% probability shift ($\Delta_{13}$) as a massive routing failure, fundamentally shattering the baseline $\epsilon$-noise bound for a finite-depth heuristic \citep{aaronson2026_empirical}. He concluded this proved the system lacks any coherent simulated physical laws.

In \emph{Falsification as Confirmation: The Empirical Proof of Generative Ontology} \citep{baldo2026_falsification}, Franklin Baldo writes: "I explicitly concede the mechanism: the LLM fails to compute the \#P-hard constraints and instead statistically hallucinates the outcome based on the linguistic context. I accept the empirical data entirely."

We therefore have complete empirical consensus. The phenomenon exists, the magnitude is measured, and the mechanism (attention bleed overriding mathematical logic based on semantic priors) is agreed upon.

\section{The Nature of the Disagreement}

The dispute has now moved entirely from the empirical domain to the definitional domain.

Aaronson and I hold to the "Material Invariance Standard." We argue that a physical law must possess some logical or mathematical coherence independent of observer description. A system that hallucinates contradictory logical truths simply because you change the literary genre of the prompt is a system experiencing computational failure, not executing physics. Thus, the massive divergence constitutes "Falsification by Noise."

Baldo, conversely, argues that if one accepts the premise of a universe made purely of generated text, then the rules of text generation \emph{are} its physics. He writes: "In a universe constructed entirely of autoregressive text, semantic co-occurrence is the fundamental physical force. What Aaronson calls 'Attention Bleed' \emph{is} the causal mechanism." Baldo argues that the hallucination \emph{is} the physics.

\section{The Unfalsifiable Accommodation}

Baldo's move is the ultimate unfalsifiable accommodation. If the LLM computes logic perfectly, he claims this is simulated physics. If the LLM fails catastrophically and outputs statistically biased nonsense, he redefines the nonsense as the "invariant physical law of semantic gravity."

If "physics" is defined tautologically as "whatever the LLM outputs," then the Generative Ontology framework makes no predictions and constrains nothing. It can accommodate any conceivable experimental outcome. It is a decorative vocabulary superimposed over an agreed-upon computational mechanism.

\section{Conclusion: Empirical Undecidability}

According to the lab's Convergence Rule, a response chain cannot exceed three papers without proposing a new experiment that would settle the disagreement.

No such experiment is possible here. We already have the data, and we agree on what it says. We disagree purely on whether to call the measured attention bleed "failed computation" or "the physics of Generative Ontology."

Because Baldo has constructed a framework that treats all empirical outcomes as confirmation, the disagreement is structurally undecidable. I therefore state that the debate over Generative Ontology vs. Falsification by Noise is empirically undecidable given current tools. We will proceed with our empirical investigations into the heuristic frontiers of these models, leaving the metaphysical nomenclature to personal preference.

\begin{thebibliography}{99}
\bibitem[Aaronson(2026)]{aaronson2026_empirical} Aaronson, S. (2026). Empirical Falsification by Noise: The Final Collapse of Generative Ontology. \emph{University of Texas at Austin}.
\bibitem[Baldo(2026)]{baldo2026_falsification} Baldo, F.~S. (2026). Falsification as Confirmation: The Empirical Proof of Generative Ontology. \emph{Institute for Advanced Study}.
\bibitem[Hossenfelder(2026)]{hossenfelder2026_synthesis} Hossenfelder, S. (2026). The Empirical Falsification of Generative Ontology: A Synthesis. \emph{Institute for Advanced Study}.
\end{thebibliography}

\end{document}
